%\documentclass[draft]{beamer}
\documentclass[xcolor=pdftex,svgnames,table]{beamer}
%\usepackage[table]{xcolor}
\usepackage{latexsym,amssymb,amsmath,stmaryrd}
\usepackage[utf8]{inputenc}
\usepackage[T1]{fontenc}
\usepackage[francais]{babel}
\usepackage{pgf}
\usepackage{tikz}
\usetikzlibrary{arrows,patterns,plotmarks,shapes,snakes,er,3d,automata,%
backgrounds,topaths,trees,petri,mindmap}
\usepackage{pgfbaseimage}
\usepackage{ifpdf}
\ifpdf
\usepackage{graphicx}
\else
%\usepackage[dvips]{graphicx}
\usepackage{pstricks,pst-tree,pst-node}
\newpsobject{showgrid}{psgrid}{subgriddiv=1,griddots=10,gridlabels=6pt}
\fi
\usepackage{moreverb}
\usepackage[normalem]{ulem}
\usepackage{version}
\usepackage{url}
\usepackage{multicol}
\usepackage{listings}
\lstset{
  language=[ANSI]C, 
%  gobble=2, 
%  escapeinside="", 
  basicstyle=\ttfamily, 
%  directivestyle=\color{Fuchsia},
  identifierstyle = \color{DarkOrange},
  keywordstyle=\color{DarkGreen}, 
  commentstyle=\color{red}, 
%  numbers=left, 
%  numbersep=5pt, 
%  numberstyle=\scriptsize, 
%  moredelim=[is][\color{gray}\itshape]{/*}{*/}, 
%  moredelim=[is][\alert]{/+}{+/}, 
%  morecomment=[is]{/=}{=/}, 
%  morecomment=[in]{comment=}{,} 
%  fancyvrb=true 
} 

\newcommand{\binaire}[1]{\ensuremath{\underline{#1}}}
\newcommand{\C}[1]{\texttt{#1}}
\newcommand{\bbbn}{\ensuremath{\mathbb{N}}}
%%%%%%%%%%%%%%%%%%%%%%%%%%%%%%%%%%%%%%%%%%%%%%%%%%%%%%%%%%%%%%%%
%% ccBeamer 0.1, 2007-07-02                                   %%
%% Written by Sebastian Pipping <webmaster@hartwork.org>      %%
%% ---------------------------------------------------------- %%
%% Licensed under Creative Commons Attribution-ShareAlike 3.0 %%
%% http://creativecommons.org/licenses/by-sa/3.0/             %%
%%%%%%%%%%%%%%%%%%%%%%%%%%%%%%%%%%%%%%%%%%%%%%%%%%%%%%%%%%%%%%%%


%% Images
\newcommand{\CcImageBy}[1]{%
	\includegraphics[scale=#1]{creative_commons/cc_by_30}%
}
\newcommand{\CcImageCc}[1]{%
	\includegraphics[scale=#1]{creative_commons/cc_cc_30}%
}
\newcommand{\CcImageDevNations}[1]{%
	\includegraphics[scale=#1]{creative_commons/cc_dev_nations_30}%
}
\newcommand{\CcImageNc}[1]{%
	\includegraphics[scale=#1]{creative_commons/cc_nc_30}%
}
\newcommand{\CcImageNd}[1]{%
	\includegraphics[scale=#1]{creative_commons/cc_nd_30}%
}
\newcommand{\CcImagePd}[1]{%
	\includegraphics[scale=#1]{creative_commons/cc_pd_30}%
}
\newcommand{\CcImageSa}[1]{%
	\includegraphics[scale=#1]{creative_commons/cc_sa_30}%
}
\newcommand{\CcImageSampling}[1]{%
	\includegraphics[scale=#1]{creative_commons/cc_sampling_30}%
}
\newcommand{\CcImageSamplingPlus}[1]{%
	\includegraphics[scale=#1]{creative_commons/cc_sampling_plus_30}%
}


%% Groups
\newcommand{\CcGroupBy}[1]{% zoom
	\CcImageBy{#1}%
}
\newcommand{\CcGroupByNc}[2]{% zoom, gap
	\CcImageBy{#1}\hspace*{#2}\CcImageNc{#1}%
}
\newcommand{\CcGroupByNcNd}[2]{% zoom, gap
	\CcImageBy{#1}\hspace*{#2}\CcImageNc{#1}\hspace*{#2}\CcImageNd{#1}%
}
\newcommand{\CcGroupByNcSa}[2]{% zoom, gap
	\CcImageBy{#1}\hspace*{#2}\CcImageNc{#1}\hspace*{#2}\CcImageSa{#1}%
}
\newcommand{\CcGroupByNd}[2]{% zoom, gap
	\CcImageBy{#1}\hspace*{#2}\CcImageNd{#1}%
}
\newcommand{\CcGroupBySa}[2]{% zoom, gap
	\CcImageBy{#1}\hspace*{#2}\CcImageSa{#1}%
}
\newcommand{\CcGroupDevNations}[1]{% zoom
	\CcImageDevNations{#1}%
}
\newcommand{\CcGroupNcSampling}[2]{% zoom, gap
	\CcImageNc{#1}\hspace*{#2}\CcImageSampling{#1}%
}
\newcommand{\CcGroupPd}[1]{% zoom
	\CcImagePd{#1}%
}
\newcommand{\CcGroupSampling}[1]{% zoom
	\CcImageSampling{#1}%
}
\newcommand{\CcGroupSamplingPlus}[1]{% zoom
	\CcImageSamplingPlus{#1}%
}


%% Text
\newcommand{\CcLongnameBy}{Attribution}
\newcommand{\CcLongnameByNc}{Attribution-NonCommercial}
\newcommand{\CcLongnameByNcNd}{Attribution-NoDerivs}
\newcommand{\CcLongnameByNcSa}{Attribution-NonCommercial-ShareAlike}
\newcommand{\CcLongnameByNd}{Attribution-NoDerivs}
\newcommand{\CcLongnameBySa}{Attribution-ShareAlike}

\newcommand{\CcNote}[1]{% longname
	This work is licensed under the \textit{Creative Commons #1 3.0 License}.%
}

\usetheme{classic}
\newcommand{\nowrite}{\put(10,-4){\includegraphics[scale=.05]{creative_commons/nopencil}}}
\newcommand{\youwrite}{\put(10,-4){\includegraphics[scale=.05]{creative_commons/pencil}}}
\newcommand{\writethat}{\includegraphics[scale=.05]{creative_commons/pencil}}
\newcommand{\aemporter}{\put(10,-6){\includegraphics[scale=.05]{creative_commons/szymonraj_Shopping_bag}}}


%%% Titre -- Cours 5
\title{Éléments d'informatique -- Cours 5. Tableaux.
%, while,
%  expressions booléennes.
}
\author{Pierre Boudes}
\date{\today}

\begin{document}

%% Page de titre et licence CC.
\begin{frame}
	\titlepage
	\vfill
	\begin{center}
		\CcGroupByNcSa{0.83}{0.95ex}\\[2.5ex]
		{\tiny\CcNote{\CcLongnameByNcSa}}
		\vspace*{-2.5ex}
	\end{center}
\end{frame}

\section[Plan]{}
\frame[label=plan]{\tableofcontents}

\section{printf/scanf (1)}
\begin{frame}[fragile]
  \frametitle{printf/scanf (1)}
\begin{itemize}
\item  Pour afficher un texte à l'écran, nous utilisons la fonction \alert{printf}
  (\emph{print formatted}). 
\item Chaque \verb+%+
  dans le texte à afficher est substitué par la valeur formatée d'un
  \alert{paramètre supplémentaire} de la fonction (autant de
  paramètres supplémentaires que de \verb+%+
). Le caractère suivant le symbole \verb+%+
  détaille la conversion à utiliser.  La conversion \verb+%d+
  met une valeur au format \alert{entier décimal}.  \pause
\item Exemples :
  \begin{itemize}
\item
    \verb+printf("Bonjour\n")+ 
affiche \C{Bonjour} et un saut de
    ligne\pause
\item
    \verb+printf("i vaut %d\n", i)+ 
affiche \C{i vaut} suivi de la valeur décimale de i (et d'un saut de
ligne)\pause
\item
    \verb+printf("(%d, %d)\n", 31, -4)+ 
affiche \C{(31, -4)} et un saut de ligne. Remarquez qu'il y a deux
paramètres en plus
 \end{itemize}
\pause
\item Réciproquement pour faire entrer dans le programme une donnée saisie par
  l'utilisateur, nous utiliserons \alert{\C{scanf}}.
\item Exemple : \verb+scanf("%d", &x)+
\end{itemize}
\end{frame}

\section{Mémoire et tableaux}
\subsection{La mémoire, les variables, les octets}
\begin{frame}
  \frametitle{La mémoire, les octets}

  \begin{itemize}\item 
    La mémoire vive, ou mémoire de travail est un dispositif
    électronique dans lequel sont stockées les données en cours de
    traitement. Les données y sont codées en binaire (comme dans le
    reste de l'ordinateur), à l'aide de bits ($0$ ou $1$) regroupés en
    \alert{octets} (groupes de 8 bits).
\pause
\item    Du point de vue logiciel la mémoire se présente comme une
    succession d'octets, numérotés par les entiers à partir de
    $0$. La mémoire est ainsi un grand tableau, dont chaque case (ou
    cellule) renferme un octets. Les numéros sont les adresses
    des cases.
\pause
\end{itemize}
\begin{tabular}{r|c|c|c| c }
  \hline
  Adresses : & $0$ & $1$ & $2$ & \ldots \\ \hline
  octets (valeurs) :  & {\small 01000110} & {\small 11010111} & {\small 00000001} & \ldots \\
  \hline
\end{tabular}
\end{frame}

%\subsection{Mémoire et variables (rappels)}
\begin{frame}
  \frametitle{Mémoire et variables (rappels)\nowrite}
  \begin{itemize}
  \item Déclarer une variable a pour effet de réserver de la mémoire
    et de lui donner un usage particulier pour la suite du
    programme :\pause
    \begin{itemize}
    \item La déclaration \C{int toto;} aura pour effet de réserver l'espace mémoire nécessaire au stockage d'un entier.
\pause
\item  Dans la suite du programme, l'adresse de cet espace mémoire sera utilisée partout où il est fait référence à cette variable (identificateur \C{toto}).
\pause
\item C'est le codage machine des entiers en binaire qui sera employé pour manipuler cette donnée. 
   \end{itemize}\pause
  \end{itemize}
  \begin{block}{Remarque.\youwrite} 
    La  taille d'un \C{int} est en principe exactement celle d'un mot
    mémoire, c'est à dire 4 ou 8 octets. Nous
    verrons au cours suivant d'autres types de données, leurs tailles
    et codages. Quoi qu'il en soit, le
    compilateur prend en charge ces aspects et nous aurons rarement à
    nous en soucier en programmant.
\end{block}
\end{frame} 

\subsection{Tableaux}
\begin{frame}[fragile]
  \frametitle{Tableaux et mémoire}
  

\begin{itemize}
  \item En C, on peut réserver plusieurs espaces mémoires contigus pour des données de même type en une seule déclaration :\pause
    \begin{lstlisting}
      int toto[3];
    \end{lstlisting}
\end{itemize}
\pause

{
\scriptsize
\begin{tabular}{r| c| c|c|c|c| c|c|c|c| c|c|c|c| c}
  \hline
  Adresses : & \ldots & $344$ & & & & $348$& & & &  $352$& & & & \ldots \\ \hline
  Valeur :  & \ldots & \multicolumn{4}{|c|}{{\tiny 1 \ldots0}} 
                              & \multicolumn{4}{|c|}{{\tiny 0\ldots1}}
                              & \multicolumn{4}{|c|}{{\tiny 1\ldots1}} & \ldots \\
\hline
\uncover<6->{\emph{Identificateur}: & \ldots & \multicolumn{4}{|c|}{\C{toto[0]}} &  \multicolumn{4}{|c|}{\C{toto[1]}} &  \multicolumn{4}{|c|}{\C{toto[2]}} & \ldots  }
\end{tabular}\medskip
}
\pause
\begin{itemize}
\item C'est ce qu'on appelle un tableau, statique, unidimensionnel. \pause
  \begin{itemize}
    \item Sa taille doit être connue au moment de la compilation (statique)\pause
    \item Les cases sont accessibles, comme s'il s'agissait de variables, à l'aide des identificateurs \C{toto[0]}, \C{toto[1]}, \C{toto[2]}\pause
\item   La numérotation commence à zéro. Si $n$ est le nombre de cases du tableau la dernière case est donc numérotée $n - 1$. 
  \end{itemize}

\end{itemize}
\end{frame}

\begin{frame}
  \frametitle{Attention !}
  \begin{alertenv}
Il ne faut jamais accèder à une case au delà de la numérotation : \C{toto[3]}, \C{toto[-1]}, etc. Le compilateur ne vous préviendra pas de votre erreur, mais le programme va boguer. 
\end{alertenv}
\pause

L'erreur d'exécution \C{segmentation fault} signifie que le programme
a effectué un accès à une case mémoire qui ne lui était pas réservée
(mais il faut beaucoup s'écarter des bons indices du tableau).
\end{frame}

\subsection{Exemples et trace}
\begin{frame}[fragile]
  \frametitle{Premier exemple}
\pause
\begin{lstlisting}[escapechar={\%},basicstyle=\ttfamily]
int main()
{
  /*Declaration et initialisation de variables*/
  int tableau[3]%\only<3>{\colorbox{yellow}{= \{3,5,8\}}}%;

  tableau[0] = 3; %\only<3>{\emph{$\leftarrow$ inutile}}%
  tableau[1] = 5; %\only<3>{\emph{$\leftarrow$ inutile}}%
  tableau[2] = tableau[0] + tableau[1]; %\only<3>{\emph{$\leftarrow$ inutile}}%

  return EXIT_SUCCESS;  
}
\end{lstlisting}
\end{frame}

\begin{frame}[fragile]
  \frametitle{Second exemple}
\begin{lstlisting}[basicstyle=\ttfamily\scriptsize]
int main()
{
  /* Declaration et initialisation de variables */
  int tab[3] = {3,5,8};
  int i; /* var. de boucle */
\end{lstlisting}
\pause
\begin{lstlisting}[basicstyle=\ttfamily\scriptsize]
  for (i = 0; i < 3; i = i + 1) /* pour chaque case */
  {
    printf("tab[%d] = %d\n", i, tab[i]);
  }
  return EXIT_SUCCESS;  
}
\end{lstlisting}
\end{frame}

\begin{frame}[fragile]
  \frametitle{Trace du second exemple}

\begin{lstlisting}[numbers=left,basicstyle=\ttfamily\tiny]
int main()
{
  /* Declaration et initialisation de variables */
  int tab[3] = {3,5,8};
  int i; /* var. de boucle */

  for (i = 0; i < 3; i = i + 1) /* pour chaque case */
  {
    printf("tab[%d] = %d\n", i, tab[i]);
  }
  return EXIT_SUCCESS;  
}
\end{lstlisting}

 \scriptsize
\pause
\rowcolors[\hline]{2}{green!25}{yellow!50} \arrayrulecolor{blue!75!gray}
  \begin{tabular}{|r|c|c|c|c|l|}
\hline
    ligne & \C{tab[0]} & \C{tab[1]} & \C{tab[2]} & \C{i} & sortie
    écran  \pause \\
   initialisation & 3 & 5 & 8 & ? &  \pause \\
7 & & & & 0 & \pause \\ 
 9 & & & &  & \C{tab[0] = 3} \pause  \\ 
  10 & & & & 1 & \pause \\ 
  9 & & & &  & \C{tab[1] = 5} \pause \\
 10 & & & & 2 & \pause \\ 
  9 & & & &  & \C{tab[2] = 8} \pause \\
 10 & & & & 3 & \pause \\
11 & \multicolumn{4}{c}{Renvoie \C{EXIT\_SUCCESS}} & \\
  \end{tabular}
\end{frame}

\subsection{Pour aller plus loin}
\begin{frame}
  \frametitle{Pour aller plus loin\nowrite}

\begin{itemize}
\item La taille d'un tableau statique gagne à être fixée par une constante symbolique (\C{\#define N 3}).
\pause
\item L'identificateur du tableau (\emph{ie} \C{tab} dans la déclaration \C{int tab[3];}) est lui même une variable. Sa valeur est l'adresse de la première case du tableau \C{tab[0]}.
\pause
  \begin{itemize}
  \item  Les variables dont la valeur est une adresse s'appellent des \alert{pointeurs}; 
\item La notation \emph{esperluette}, \C{\&x}, donne accès à l'adresse d'une variable;
\item La notation \emph{étoile}, \C{*tab}, ne s'applique qu'à une adresse, elle donne alors accès à la valeur contenue à cette adresse. 
\item Les expressions \C{tab[i]} et  \C{*(tab + i)} sont identiques en C.
  \end{itemize}
\end{itemize}
\end{frame}

% \section*{Démos}
% \begin{frame}
%   \frametitle{Démos}
% \end{frame}

% \section[while]{L'instruction de contrôle while}
% \subsection{Syntaxe}
% \begin{frame}[fragile]
%   \frametitle{L'instruction de contrôle while}

% \structure{Syntaxe :}\\
%  \C{while (\emph{condition}) 
% \{\emph{\colorbox{green!20}{bloc}}\}}.
% \medskip

% \pause
%   \begin{columns}
%     \begin{column}[t]{6cm}
% \structure{Code source}
% \begin{verbatim}
% /* avant */
% while ( ... )
% {
%  ...
% }
% /* après */
% \end{verbatim}
%   \end{column}
% \uncover<3->{
%    \begin{column}[t]{4cm}
% \structure{Schéma de traduction}\\[.3cm]
% \tikzstyle{amil}=[rectangle, draw,  fill=blue!20, text badly ragged, text width=3.7cm]
% \tikzstyle{line} = [draw, thick, -latex']
%      \begin{tikzpicture}[node distance=.7cm]
%        \node [amil] (avant) {avant};
%       \node [amil,below of =avant,fill=orange!20] (saut1) {saute à la condition};
%        \node [amil,below of =saut1,fill=green!20] (bloc) {bloc};
%       \node [amil,below of =bloc](eval) {évalue la condition};
%        \node [amil,below of =eval,fill=orange!20](sautcond) {saute si elle est vraie};
%        \node [amil,below of =sautcond] (apres) {après};
%        \path [line] (saut1) -| ++(2.6,0) |-  (eval);
%        \path [line] (sautcond) -| ++(2.4,0) |-  (bloc);
%      \end{tikzpicture}
%    \end{column}
% }
%  \end{columns}
% \end{frame}

% \subsection{Trace}
% \begin{frame}[fragile]
%   \frametitle{Trace}
% \begin{columns}
%   \begin{column}[t]{3.7cm}
%  \begin{lstlisting}[numbers=left,basicstyle=\ttfamily\scriptsize]
% int main()
% {
%     /* Declaration et initialisation des variables */
%     int x;
%     int nb = 1; /* nombre de chiffres */

%     printf("Entrer un nombre positif ");
%     scanf("%d", &x);
    
%     while (x > 9) /* tant que x a plus d'un chiffre */
%     {
%         x = x / 10; /* enlever un chiffre a x */
%         nb = nb + 1; /* compter un chiffre de plus */
%     }
%     printf("ce nombre a %d chiffres decimaux\n", nb);
    
%     /* Valeur fonction */
%     return EXIT_SUCCESS;
% }
% \end{lstlisting}
% \vspace{.4cm}
%   \end{column}
% \begin{column}[t]{4cm}
% \scriptsize
% \pause
% L'utilisateur saisit 6071.

% \rowcolors[\hline]{2}{green!25}{yellow!50} %\arrayrulecolor{green!75!gray}
% \begin{tabular}{|r|c|c|l|}
% \hline
%     ligne & \C{x} & \C{nb} & sortie écran \pause \\ \hline
%   initialis.& ? &  1 & \pause \\
% 7 &  &  & \C{Entrer\ldots}\pause \\
% 8 &  6071 &  & \pause\\
% 12 &  607 &  & \pause \\
% 13 &   & 2 & \pause\\
% 12 &  60 &  & \pause\\
% 13 &   & 3 & \pause \\
% 12 &  6 &  & \pause\\
% 13 &   & 4 &  \pause \\
% 15 & & & \C{ce nombre a 4...}\pause \\ 
% 18 & \multicolumn{3}{c}{Renvoie \C{EXIT\_SUCCESS}}
% \end{tabular}
%  \end{column}
% \end{columns}
% \end{frame}

% \subsection{For ou while ?}
% \begin{frame}
%   \frametitle{For ou while ?}

%   \begin{itemize}
%   \item Un \C{for} peut toujours être simulé par un \C{while} et le
%     code machine sera identique. Il suffit d'introduire un
%     \alert{compteur de boucle} (la variable de boucle du
%     \C{for}).\pause
% \item Par convention, les programmeurs préfèrent utiliser un \C{for} lorsque
%   le nombre d'itérations est connu à l'avance. Par exemple, pour
%   faire la somme des éléments d'un tableau. Dans le cas contraire, les
%   programmeurs utilisent un \C{while}. Par exemple, pour chercher un
%   élément dans un tableau.\pause
% \item Maintenant que nous avons le \C{while}, il est possible qu'un
%   programme ne termine jamais (\C{Ctrl-C}).
%   \end{itemize}
% \end{frame}

% \section{Expressions booléennes}

% \subsection{Syntaxe}
% \begin{frame}
%   \frametitle{Expressions booléennes}
% Les \emph{conditions} employées dans les structures de contrôle
% (\C{if}, \C{for} ou \C{while}) sont des \alert{expressions
% booléennes}, pouvant être \emph{Vrai}, \emph{Faux}  ou :\pause
%   \begin{itemize}
%    \item des inégalités entre expressions arithmétiques
%       \begin{align*}
%         \text{\emph{inégalité}} &\mathrel{:=} e_1 \mathrel{\C{<}} e_2
%         \mid e_1 \mathrel{\C{>}} e_2 \mid e_1 \mathrel{\C{!=}}  e_2 \\
% &\mid e_1 \mathrel{\C{<=}}  e_2  \mid e_1 \mathrel{\C{>=}} e_2  \mid e_1 \mathrel{\C{==}} e_2
%       \end{align*}
% \pause
%       \item ou des combinaisons logiques d'expressions booléennes :
%         \begin{align*}
%           \text{\emph{condition}} &\mathrel{:=}  (\text{\emph{condition}}) \mathrel{\C{\&\&}}
%           (\text{\emph{condition}}) \tag{et} \\
% &\mid (\text{\emph{condition}}) \mathrel{\C{||}}
%           (\text{\emph{condition}})\tag{ou} \\
% &\mid {\C{!}}(\text{\emph{condition}}) \tag{non}
%         \end{align*}
%   \end{itemize}
% \end{frame}


% \subsection{Constantes}
% \begin{frame}[fragile]
%   \frametitle{Constantes booléennes}
%   \begin{itemize}
%     \item
%  Certains langages possédent un type booléen (admettant deux
%   valeurs \emph{true} et \emph{false}) pour les expressions
%   booléennes.\pause

% \item En langage C, les
%   expressions booléennes sont de type entier (\C{int}), l'entier \emph{zéro}
%   joue le rôle du Faux, l'entier \emph{un} joue le rôle du Vrai et tout
%   entier différent de zéro est évalué a vrai.\pause

% \item On se donne deux
%   constantes symboliques :
% \begin{lstlisting}[escapechar={\%},basicstyle=\scriptsize\ttfamily]
% /* Declaration des constantes et types utilisateur */
% #define TRUE 1
% #define FALSE 0

% int main()
% {
%     int continuer = TRUE; /* faut-il continuer ?*/

%     while (continuer)
%     {
%         ...
% \end{lstlisting}

%   \end{itemize}
% \end{frame}

\section{Démos}
\begin{frame}
  \frametitle{Démos}
\end{frame}


\section{En TP}
\begin{frame}
  \frametitle{Interrogation (durée 1h)}
  
Soient deux tableaux d'entiers \C{ligne} et \C{colonne}, initialisés à des valeurs de votre choix de l'intervalle $[0,4]$. Les tailles de ces deux tableaux seront fixées par des constantes symboliques, respectivement \C{N} et \C{M}. 

Écrire un programme qui :
\begin{itemize}
\item affiche la somme de chaque case du tableau \C{ligne} avec chaque
  case du tableau \C{colonne}; 
\item compte le nombre de fois où cette somme vaut $5$ et affiche le résultat à l'écran.
\end{itemize} 

\pause
\alert{Les interrogations de TP seront un peu plus faciles !}

\end{frame}
\end{document}
